\section{Results}
\label{sec:results}

% NOTE: no fresh numbers exist yet (you're feeding those in later), and I
% deliberately did NOT reuse any figure from IMPLEMENTATION.md's "Analysis of
% tendency" section, since you flagged that file as an obsolete pre-D=9
% version. What follows is the narrative scaffold + placeholders your outline
% asks for, plus one thing I could derive directly from the math/code rather
% than from a specific run (the lambda=0 remark) -- that part I *am*
% confident in, independent of which run's numbers eventually go in.

\subsection{Data and candidate pool}
\label{sec:results:data}
We use every SNOMED~CT concept mapped by our institution's semantician team
in HUGData as the mapped set $M$. \notetoself{"HUGData" is used here exactly
as it appears in your own \texttt{config.py} path names -- flag if this
should instead be spelled out or anonymized for the paper.} Candidate tokens
are restricted to the $D$-hop neighborhood of $M$ (Section~\ref{sec:method:candidates}),
which narrows the pool from $|V| \approx 400{,}000$ SNOMED concepts to
\needresults{current candidate-pool size under \texttt{max\_dist\_candidate = 9}}.
\doubt{Your outline states this narrowing as "400k -> 70k," which matches the
75{,}634-node figure from the old (obsolete) $D=3$ analysis -- not the current
$D=9$ setting. A larger $D$ grows the ego-graph radius substantially, so the
current pool is very likely much larger than 70k, possibly a large fraction of
all of $V$. Please recompute this from a current run of
\texttt{1.graph\_prepare.ipynb} before this number goes in the paper.}

\subsection{General trend across $k$}
\label{sec:results:trend}
\needresults{plot(s) of each metric from Section~\ref{sec:method:metrics}
against $k$, one line per candidate-selection method, at a fixed reporting
$\lambda$ (presumably $\lambda = 1$, consistent with
Section~\ref{sec:method:metrics}'s evaluation convention). Send me the
underlying table/plot and I will write the narrative text around it (not
just "Figure~X shows Y" but why the shape looks that way given the method).}

\subsection{Comparison across baseline families}
\label{sec:results:families}
For each family from Section~\ref{sec:method:baselines} -- random,
degree-based, centrality-based, and ours (with its hard-greedy ablation) --
we plot the upper envelope of metric performance across the family's members,
so that the comparison is against each family's \emph{best} member rather
than diluted by its weaker ones.

\needresults{once the sweep is in, this subsection should state, per family,
which member forms the envelope at which $k$, and characterize \emph{why}.
Below are hypotheses that follow from the method design in
Section~\ref{sec:method}, phrased so they're easy to confirm or reject
against the actual numbers rather than asserted as findings:}
\begin{itemize}
\item \doubt{Degree-based methods (especially \emph{most children}, which
only looks at the \texttt{IS\_A} subgraph) may tend to select candidates
whose neighborhoods overlap heavily -- e.g.\ several siblings under the same
popular parent -- which would show up as lower discriminative-power metrics
(uniqueness entropy / unique rate) than a coverage-aware method at the same
$k$, even where raw semantic coverage is comparable. Worth checking directly
rather than assuming.}
\item \doubt{Centrality-based methods (PageRank family, closeness,
eigenvector) rank by global random-walk or shortest-path structure rather
than by $M$-specific coverage; personalized PageRank is the one exception
(it restarts on $M$). This predicts personalized PageRank should track our
method more closely than plain PageRank or closeness does -- again, a
hypothesis to check, not yet a result.}
\item \doubt{eigenvector centrality's DAG degeneracy (flagged in
Section~\ref{sec:method:baselines}) predicts its envelope contribution should
flatten out once $k$ passes whatever fraction of the pool actually carries
non-negligible score -- if that's what the current run shows, it's worth
stating explicitly as an explanation rather than leaving the flat curve
unexplained.}
\end{itemize}

\subsection{Effect of the distance penalty \texorpdfstring{$\lambda$}{lambda}}
\label{sec:results:lambda}
\textbf{The degenerate case $\lambda = 0$.} This case can be characterized
exactly from Eq.~\eqref{eq:coverage} itself, independent of any particular
run: at $\lambda = 0$, every term of the "otherwise" branch vanishes, so
$S_h(u,T) = S_0(u,T) = \mathbf 1[u \in T]$ for \emph{every} horizon $h$, not
only $h=0$ -- coverage never propagates. Substituting into $F_D$,
\[
F_D(T)\big|_{\lambda=0} = \frac{1}{|M|}\sum_{c \in M} \mathbf 1[c \in T] = \text{exact\_rate}(M, T).
\]
In other words, selecting tokens at $\lambda = 0$ is not "coverage selection
with a poor setting" but optimization of a \emph{different} objective
entirely: how many mapped concepts the selector picked as tokens of
themselves. Any candidate not itself in $M$ has zero marginal gain under this
objective, so the selector has no signal to rank non-$M$ candidates at all,
and its behavior past the point where $M$ is exhausted is effectively
arbitrary tie-breaking. \notetoself{This derivation only needs
Eq.~\eqref{eq:coverage} and the definition of exact\_rate
(Section~\ref{sec:method:metrics}) -- it holds regardless of which run's
numbers eventually back it up, so I'm confident stating it as a proposition
rather than an empirical observation. \needresults{the actual $\lambda=0$
curve, to show this pathology concretely rather than just algebraically.}}

\textbf{Trade-off across $\lambda \in (0, 1]$.}
\needresults{once available: how each metric family from
Section~\ref{sec:method:metrics} moves as $\lambda$ increases from $0$ toward
$1$. A reasonable prior, to be confirmed or corrected by the actual sweep, is
that metrics tied directly to reaching \emph{any} token (semantic coverage,
unk rate, uncovered rate) should improve monotonically as $\lambda \to 1$
(evaluation is always at $\lambda=1$, so selection and evaluation become
better aligned), while metrics sensitive to \emph{how far} a token was found
(distance score, conciseness) may instead peak at some intermediate
$\lambda$, since $\lambda=1$ selection has no preference for nearby tokens
over merely reachable ones. \doubt{state as hypothesis or confirmed finding
depending on what the fresh sweep shows -- do not carry over the specific
numbers from the old, obsolete analysis even if the qualitative shape turns
out to match.}}

\textbf{Token reuse rate.} \doubt{your outline flags this metric as "need to
be implemented" -- confirmed: \texttt{eval.py} does not currently compute
anything under this name. This subsection stays a placeholder until that
metric exists in code; happy to help define/implement it once you specify
what "reuse" should mean here (e.g.\ how many distinct mapped concepts a
single token ends up representing, across the whole context-tree
population).}

\textbf{Practical takeaway.} \doubt{your outline's closing point, "best
config depends on personal choice" -- worth turning into one concrete
sentence once the trade-off curves exist, e.g.\ naming which $\lambda$ range
is a reasonable default and what a practitioner would trade away by moving
off it in either direction.}

\notetoself{Your PDF's own Section 7 ("Experiments") also proposes sweeping
$D$ itself -- $D \in \{1,2,3,4,5\}$, measuring $|V_D|$ growth, reachable
relationship counts, and degree distributions -- as a separate axis from $k$
and $\lambda$. That sweep doesn't appear anywhere in this Results outline,
which otherwise treats $D$ as fixed. Worth deciding explicitly whether a
$D$-sensitivity analysis belongs in this paper's Results or is out of scope.}
